\documentclass[10pt]{beamer}

% --- Modern Theme ---
\usetheme{metropolis}
\usecolortheme{default}
\setbeamercolor{background canvas}{bg=white}

% --- Packages ---
\usepackage[utf8]{inputenc}
\usepackage{booktabs}
\usepackage{array}
\usepackage{tabularx}
\usepackage{amsmath}
\usepackage{amssymb}
\usepackage[scale=2]{ccicons}
\usepackage{pgfplots}
\usepgfplotslibrary{dateplot}
\usepackage{xspace}

% --- Title Information ---
\title{Foundations of Artificial Intelligence}
\subtitle{Part II: Problem Solving and Search}
\institute{Department of Computer Science}
\date{}

\begin{document}

\maketitle

\begin{frame}{Lecture Roadmap}
    \setbeamertemplate{section in toc}[sections numbered]
    \tableofcontents[hideallsubsections]
\end{frame}

\begin{frame}{Learning Objectives}
    By the end of this lecture, you should be able to:
    \begin{itemize}
        \item Formulate an AI task as a search problem over a state space.
        \item Compare uninformed strategies: BFS, DFS, and Uniform-Cost Search.
        \item Explain informed search with Greedy Best-First and A*.
        \item Use heuristic properties (admissibility and consistency) to reason about optimality.
        \item Choose a strategy based on completeness, optimality, time, and space trade-offs.
    \end{itemize}
\end{frame}

\section{From Problem Formulation to Search}

\begin{frame}{Why Search Matters in AI}
    Problem-solving agents follow a simple loop:
    \begin{center}
        \fbox{Formulate $\rightarrow$ Search $\rightarrow$ Execute}
    \end{center}
    \vspace{0.3cm}
    \begin{itemize}
        \item \textbf{Formulate:} Convert the task into states, actions, and goals.
        \item \textbf{Search:} Find an action sequence from initial state to a goal state.
        \item \textbf{Execute:} Apply the sequence in the environment.
    \end{itemize}
    \vspace{0.2cm}
    \textit{Core idea from the book: intelligence emerges from representation + algorithm.}
\end{frame}

\begin{frame}{State Space Representation}
    A search problem is specified as:
    \[
        \langle S, A, s_0, G, c \rangle
    \]
    \begin{itemize}
        \item $S$: set of states.
        \item $A(s)$: actions available in state $s$.
        \item $s_0$: initial state.
        \item $G$: goal test or set of goal states.
        \item $c(s,a,s')$: step cost function.
    \end{itemize}
    \vspace{0.2cm}
    A solution is a path from $s_0$ to a goal; an \textbf{optimal} solution has minimum total path cost.
\end{frame}

\begin{frame}{Search Tree vs State Graph}
    \begin{columns}[T,onlytextwidth]
        \column{0.48\textwidth}
            \alert{Search Tree View}
            \begin{itemize}
                \item Nodes are paths.
                \item Same state may appear many times.
                \item Can become infinite with loops.
            \end{itemize}
        \column{0.48\textwidth}
            \alert{State Graph View}
            \begin{itemize}
                \item Nodes are unique states.
                \item Graph-search tracks visited states.
                \item Usually reduces repeated work.
            \end{itemize}
    \end{columns}
    \vspace{0.3cm}
    \textbf{Practical rule:} Prefer graph search unless memory constraints are severe.
\end{frame}

\begin{frame}{Running Example: Route Finding}
    \small
    \begin{tabular}{p{2.4cm}p{7.9cm}}
        \toprule
        \textbf{Component} & \textbf{Example} \\
        \midrule
        States & Cities on a map \\
        Actions & Drive to an adjacent city \\
        Initial state & Milan \\
        Goal test & Reach Rome \\
        Step cost & Road distance in km \\
        \bottomrule
    \end{tabular}
    \vspace{0.2cm}


    This classic setup lets us compare BFS, DFS, Uniform-Cost, Greedy, and A* on the same problem.
\end{frame}

\section{Uninformed Search}

\begin{frame}{Uninformed Search: What It Knows}
    Uninformed (blind) search uses only:
    \begin{itemize}
        \item successor function,
        \item goal test,
        \item path cost definition (if relevant).
    \end{itemize}
    \vspace{0.2cm}
    It has no estimate of how close a state is to the goal.

    \vspace{0.2cm}
    \textbf{Main strategies in this lecture:}
    \begin{itemize}
        \item Breadth-First Search (BFS)
        \item Depth-First Search (DFS)
        \item Uniform-Cost Search (UCS)
    \end{itemize}
\end{frame}

\begin{frame}{Breadth-First Search (BFS)}
    \begin{itemize}
        \item Expands shallowest frontier node first (FIFO queue).
        \item Complete if branching factor is finite.
        \item Optimal when all step costs are equal.
    \end{itemize}
    \vspace{0.2cm}
    Typical complexity (tree view):
    \[
        \text{Time } = O(b^d), \quad \text{Space } = O(b^d)
    \]
    where $b$ is branching factor and $d$ is depth of shallowest goal.

    \vspace{0.2cm}
    \alert{Strength:} Reliable shallow solution.
    \hfill
    \alert{Weakness:} Very high memory use.
\end{frame}

\begin{frame}{Depth-First Search (DFS)}
    \begin{itemize}
        \item Expands deepest frontier node first (stack or recursion).
        \item Space-efficient because it stores one path plus siblings.
        \item Not optimal; can be incomplete in infinite-depth spaces.
    \end{itemize}
    \vspace{0.2cm}
    Typical complexity (tree view):
    \[
        \text{Time } = O(b^m), \quad \text{Space } = O(bm)
    \]
    where $m$ is maximum depth.

    \vspace{0.2cm}
    \alert{Strength:} Low memory footprint.
    \hfill
    \alert{Weakness:} Can go deep in wrong direction.
\end{frame}

\begin{frame}{Uniform-Cost Search (UCS)}
    \begin{itemize}
        \item Expands node with smallest path cost $g(n)$ first.
        \item Implemented with a priority queue.
        \item Equivalent to Dijkstra-style expansion on positive-cost edges.
    \end{itemize}
    \vspace{0.2cm}
    Guarantees:
    \begin{itemize}
        \item Complete if each step cost is at least some $\epsilon > 0$.
        \item Optimal for general nonnegative step costs.
    \end{itemize}
    \vspace{0.2cm}
    \textit{Key difference from BFS: UCS optimizes path cost, not depth.}
\end{frame}

\begin{frame}{Uninformed Strategy Comparison}
    \footnotesize
    \setlength{\tabcolsep}{4pt}
    \renewcommand{\arraystretch}{1.12}
    \begin{tabularx}{\textwidth}{>{\raggedright\arraybackslash}p{2.1cm} *{4}{>{\raggedright\arraybackslash}X}}
        \toprule
        \textbf{Strategy} & \textbf{Complete?} & \textbf{Optimal?} & \textbf{Time} & \textbf{Space} \\
        \midrule
        BFS & Yes (finite $b$) & Yes (equal costs) & Exponential in $d$ & Exponential in $d$ \\
        DFS & No (infinite depth) & No & Exponential in $m$ & Linear in depth \\
        UCS & Yes ($c \geq \epsilon$) & Yes & Exponential in $C^*/\epsilon$ & Similar to time \\
        \bottomrule
    \end{tabularx}
\end{frame}

\section{Informed Search}

\begin{frame}{From Blind to Informed Search}
    Informed search adds a heuristic estimate:
    \[
        h(n) \approx \text{estimated cost from } n \text{ to a goal}
    \]
    \vspace{0.2cm}
    Better estimates usually mean fewer expansions.

    \vspace{0.2cm}
    Two central strategies:
    \begin{itemize}
        \item \textbf{Greedy Best-First:} expand smallest $h(n)$.
        \item \textbf{A*:} expand smallest $f(n)=g(n)+h(n)$.
    \end{itemize}
\end{frame}

\begin{frame}{Greedy Best-First Search}
    \begin{itemize}
        \item Chooses node that appears closest to goal by $h(n)$.
        \item Often very fast with good heuristics.
        \item Can be misled into long or costly detours.
    \end{itemize}
    \vspace{0.2cm}
    Evaluation:
    \[
        f(n)=h(n)
    \]
    \vspace{0.2cm}
    \alert{Not generally optimal.}
    \hfill
    \alert{Not always complete in graph variants without safeguards.}
\end{frame}

\begin{frame}{A* Search}
    A* balances cost-so-far and estimated cost-to-go:
    \[
        f(n)=g(n)+h(n)
    \]
    \begin{itemize}
        \item $g(n)$: exact path cost from start to $n$.
        \item $h(n)$: heuristic estimate from $n$ to goal.
    \end{itemize}
    \vspace{0.2cm}
    Intuition:
    \begin{itemize}
        \item Greedy alone over-trusts $h$.
        \item UCS alone ignores goal direction.
        \item A* combines both signals.
    \end{itemize}
\end{frame}

\section{Heuristics and Optimality}

\begin{frame}{Heuristic Quality Conditions}
    \textbf{Admissible heuristic:}
    \[
        0 \le h(n) \le h^*(n)
    \]
    It never overestimates true remaining cost $h^*(n)$.

    \vspace{0.25cm}
    \textbf{Consistent heuristic (monotone):}
    \[
        h(n) \le c(n,a,n') + h(n')
    \]
    for every edge $(n,a,n')$.

    \vspace{0.25cm}
    Consistency implies admissibility (with $h(\text{goal})=0$) and gives nondecreasing $f$-values along paths.
\end{frame}

\begin{frame}{When Is A* Optimal?}
    For tree search, A* is optimal if $h$ is admissible.

    \vspace{0.25cm}
    For graph search, A* is optimal with a consistent heuristic and proper closed-list handling.

    \vspace{0.25cm}
    Why (high-level):
    \begin{enumerate}
        \item A* always expands the node with smallest $f$.
        \item For admissible $h$, $f(n)$ never overestimates best solution through $n$.
        \item Therefore, the first goal removed from frontier has cost $C^*$.
    \end{enumerate}
\end{frame}

\begin{frame}{Designing Heuristics in Practice}
    Reliable methods:
    \begin{itemize}
        \item \textbf{Relaxed problems:} remove constraints and solve easier version.
        \item \textbf{Abstractions:} compress state description while preserving lower bounds.
        \item \textbf{Pattern databases:} precompute exact costs for abstract subproblems.
    \end{itemize}
    \vspace{0.2cm}
    Quality trade-off:
    \[
        \text{more accurate } h \Rightarrow \text{fewer expansions, often more computation per node}
    \]
\end{frame}

\begin{frame}{Example Heuristic: Straight-Line Distance}
    In route planning, let $h(n)$ be straight-line distance to destination.
    \begin{itemize}
        \item Usually admissible: roads cannot be shorter than straight-line geometry.
        \item Often consistent due to triangle inequality.
    \end{itemize}
    \vspace{0.2cm}
    Result:
    \begin{itemize}
        \item Greedy becomes fast but may pick costly routes.
        \item A* with this heuristic is both efficient and optimal (under assumptions above).
    \end{itemize}
\end{frame}

\begin{frame}{Common Pitfalls and Fixes}
    \begin{itemize}
        \item \textbf{Pitfall:} Confusing state with node/path record.
        \item \textbf{Fix:} Track parent, action, and path cost per node explicitly.

        \item \textbf{Pitfall:} Using heuristic that overestimates and expecting optimality.
        \item \textbf{Fix:} Validate admissibility/consistency on representative transitions.

        \item \textbf{Pitfall:} Ignoring repeated states.
        \item \textbf{Fix:} Use graph-search with reached/closed set and best-cost updates.
    \end{itemize}
\end{frame}

\begin{frame}{From Search to Algorithmic Intelligence}
    Why this topic is foundational:
    \begin{itemize}
        \item You model the world as formal states and transitions.
        \item You reason about guarantees: completeness and optimality.
        \item You make principled trade-offs: speed, memory, and solution quality.
        \item You design heuristics that inject domain knowledge into generic algorithms.
    \end{itemize}
    \vspace{0.2cm}
    \textit{This is where intelligent behavior becomes an engineering discipline.}
\end{frame}

\begin{frame}{Key Takeaways}
    \begin{enumerate}
        \item State-space representation determines what search can do.
        \item BFS, DFS, and UCS differ mainly in frontier ordering and guarantees.
        \item Greedy is goal-directed but not reliably optimal.
        \item A* is optimal with admissible/consistent heuristics and correct graph handling.
        \item Better heuristics can drastically reduce search effort.
    \end{enumerate}
\end{frame}

\begin{frame}{Reference Reading for This Part}
    \begin{itemize}
        \item \textbf{Russell \& Norvig (3rd/4th ed.):} Problem-solving and search chapters.
        \item Focus sections: problem formulation, tree/graph search, uninformed and informed search, heuristic properties.
    \end{itemize}
    \vspace{0.25cm}
    Suggested self-check:
    \begin{enumerate}
        \item Can you write $\langle S, A, s_0, G, c \rangle$ for a new domain?
        \item Can you justify whether a heuristic is admissible?
        \item Can you predict whether an algorithm is complete/optimal in that domain?
    \end{enumerate}
\end{frame}

\begin{frame}[standout]
    Questions?
\end{frame}

\end{document}
